\subsection{Assigning packets to plates}\label{w:assignment}

The plate covering is used only to organize the original compact
packets.  It does not change their profiles.  Assign a label to a plate
if its central line comes within distance $X^{C_8}$ of that plate
somewhere in the enlarged tube-coordinate domain and its normal
velocity deviation for the same plate is at most $X^{C_9}/L$.
If there are several choices, choose one.  Thus the assigned coefficient
energies are disjoint.  Increase each assigned plate's width by a fixed
power of $X$ so that it contains the compact packet support in the
domain: the normal displacement over time $O(L)$ is at most $X^{C_9}$,
in addition to the initial distance and the fixed packet width.

Every unassigned label has moment on $\mu|_{Y_0}$ at most
$mL^{n+1}/D_1$, after choosing the powers in order.  Indeed choose
$C_8$ larger than the plate-width power and truncate its moment at
distance $X^{C_8}/2$, paying the tail by the same annular estimate as
in \Cref{w:heavy-reduction}.  If its line is ever this close to a
plate but is unassigned, its normal velocity exceeds $X^{C_9}/L$.
The time spent within that distance is at most
\[
 CL\frac{X^{C_8}+X^{C_{\mathrm{plate}}}}{X^{C_9}}
\]
for each plate, in tube coordinates.  Cover those intervals by width
boxes and sum over the polynomial number of plates.  Choosing $C_9$
large makes the resulting moment smaller than the asserted bound.
Inner induction at $D_1$ therefore removes the unassigned field on
only a small fraction of $Y_0$.  The assigned sum is at least a fixed
multiple of $\lambda$ on a set of mass comparable to $M_Y$.

Choose a dyadic inverse cap width $h$, nested with $R$, such that
\begin{equation}\label{w:eq-h}
 b_1=L/h\asymp X^{C_{10}}.
\end{equation}
The exponent $C_{10}$ is larger than the preceding plate-count and
width powers and the loss in \eqref{w:eq-incidence-mass}.  Intersections
of two enlarged plates whose projective normals differ by at least
$1/h$ have small total mass.  In tube coordinates such an intersection
fits in a box with $n-1$ lengths $O(L)$ and two remaining widths at
most $X^Ch$.  Covering by balls of the latter radius gives mass at most
$X^CL^{n-1}h$.  The same bound, with a larger fixed power, pays all
pairs of plates.  By \eqref{w:eq-incidence-mass} and
\eqref{w:eq-h}, this is negligible relative to $N_Y$.

Group the plates by a bounded-overlap grid of projective normals at
resolution $1/h$, choosing labels disjointly.  At a retained point only
boundedly many clusters contribute, so one cluster field has modulus
$\gtrsim\lambda$.  Cluster energies partition a subenergy of the
original array.  Within each cluster all frequencies lie within $C/h$
of a single affine frequency hyperplane.  To see this, the assignment
condition gives an almost tangent direction of bounded velocity for
each plate, so the spatial part of its normal is bounded below; changing
the normal by $O(1/h)$ changes its tangent-frequency hyperplane by the
same order.  A bounded spatial rotation, shear and modulation now put
this hyperplane in tangential frequency coordinates of dimension
$k=n-1$.  These transformations are used only in the following
nonrecursive estimates.

For each cluster, subdivide labels in the original frequency axes into
$1/R$-bins and then $1/h$-bins.  On a macro cube where the cluster is
high, call a coarse bin significant if its field has a witness there
of modulus at least $cR^{-n}\lambda$.  The nonsignificant fields sum
to at most a small multiple of $\lambda$ at a tested point.

Except for a small total mass, the tangential projections of the
significant centers cannot be covered by $O(C_{11}\log X)$ strips of
width $C_*/R$ about affine $(k-1)$-planes.  Here $C_*$ is any sufficiently
large fixed dimensional constant, and $C_{11}$ is the later fixed band
cutoff exponent.  Such a cover, together with the cluster's $C/h$
normal width, would put the original bins in codimension-two strips
in full frequency space.  Apply \Cref{w:angular} to the cluster subarray
on each adverse macro.  Summing disjoint energies and using $p/2>1$
preserves the bound.  Increasing $B$ absorbs the fixed logarithmic
power.  All ensuing coarse-bin witnesses are allowed to be anywhere
within $O(R^4)$ of their tested points.

\subsection{Flattened packets and their two energy budgets}

For a fine bin $\theta$ centered at $v$, use the original-axis map
$F_\theta=A_{h,v}$.  After removing its carrier, the fine-bin field is
\[
 h^{-n/2}U'_\theta(F_\theta z),
\]
where $U'_\theta$ is a child array of size $b_1$.  Partition its native
spacetime into smooth compact unit cells $T$, with a fixed enlargement
$T^*$ containing each cutoff support.  The resulting packets have
carrier $v$, width $h$ and duration $h^2$.  Their sizes and associated
densities can be taken as
\begin{align}
 s_{\theta,T}&=Ch^{-n/2}\max_{|\alpha|\leq N_*}
                     \sup_{T^*}|\partial^\alpha U'_\theta|,
                     \label{w:eq-flatsize}\\
 d_{\theta,T}&=(mh^{n+1})^{-1}
  \int(1+\dist(F_\theta z,T))^{-P_*}\,\dd\mu(z).
                     \label{w:eq-flatdensity}
\end{align}
Here $P_*,N_*$ are the fixed orders committed to in step
\textup{(ii)} of \Cref{w:choices}.  The envelope order $P=P_*$ was
chosen strictly above both $M+n+2$ and the lower envelope order in
\Cref{f:frame}.  Use it both for the full square-envelope weight
$w_T$ and for the density kernel.  Thus
\[
 \int w_T(z)\,\dd\mu(z)=mh^{n+1}d_{\theta,T},
\]
before replacing densities by their dyadic band values.  The order
$N_*$ includes the frame derivative order and the fixed reserves for the
normal Fourier expansion and cell Sobolev estimates.  Dividing each
localized profile by $s_{\theta,T}$ then gives the same fixed
low-order profile bound in every frame application.  Extra original
regularity needed by the sparse estimate does not redefine these
sizes or frame inputs.

Keep the cells needed on the physical domain and its macro
enlargements.  The pushed measure has ball parameter $Cmh^{n+1}$, so
$d_{\theta,T}\lesssim1$.

\begin{lemma}[Energy budgets after flattening]\label{w:flatten}
For each physical time bin of length $h^2$, including its fixed
enlargement,
\begin{equation}\label{w:eq-flat-unweighted}
 \sum_{\mathrm{bin}}s_{\theta,T}^2\lesssim h^{-n}.
\end{equation}
There is also the global bound, summed over all clusters, fine bins,
cells and time bins with bounded overlap,
\begin{equation}\label{w:eq-flat-weighted}
 \sum_{\theta,T}d_{\theta,T}^{\beta}s_{\theta,T}^2
 \lesssim_B h^{-n}b_1^{2n/(n+1)}D^{-c_s}.
\end{equation}
The first assertion is only a per-time-bin statement.
\end{lemma}

\begin{proof}
A physical time interval of length $h^2$ has bounded length in the
child time variable, even though a full child packet has duration
$b_1^2$.  At each such time, the Gram matrix of child atoms, also after
the fixed derivatives in \eqref{w:eq-flatsize}, has summable rows.
Integration by parts gives decay in the scaled carrier separation
$b_1(\omega-\omega')$, while compact spatial profiles localize
intercepts.  Transported phase-space counting bounds the number of
labels at each scaled distance.  All child frequencies are in a fixed
bounded box, so differentiation introduces no positive power of $b_1$.
Consequently
\[
 \sum_\theta\norm{\partial^\alpha U'_\theta(\cdot,t)}_2^2
 \lesssim\sum_\theta E_\theta\leq1
\]
for the required derivatives.  Unit-cell Sobolev, summed in space and
over the bounded child-time interval, proves
\eqref{w:eq-flat-unweighted}.

For the weighted bound apply \Cref{s:sparse}, including its fixed
derivative versions, to every full child array.  The pushforward of
$\mu$ has ball parameter $Cmh^{n+1}$ and its packet moment at scale
$b_1$ retains the same $D$, since
$h^{n+1}b_1^{n+1}=L^{n+1}$.  The choice
$b_1\asymp(BD)^{C_{10}}$ ensures that $D$ is at least a fixed positive
power of $b_1$ once $D$ is larger than a constant depending on $B$;
the remaining bounded $b_1,D$ cases cost only a constant depending on
$B,C_{10}$.  Also $D\leq b_1$ by the choice of $C_{10}$.

Convolve the pushed measure by the isotropic polynomial kernel used
in \eqref{w:eq-flatdensity}, and restrict to the suitably enlarged
time range.  Its ball and packet-moment bounds are preserved: translation
of a packet weight costs at most its fixed polynomial order, and
$P_*>M+n+2$ makes that cost integrable.  The convolved measure has
mass at least a constant times $mh^{n+1}d_{\theta,T}$ on the cells
where the suprema are taken.  The sparse estimate therefore gives
\eqref{w:eq-flat-weighted}, after division by the normalization in
\eqref{w:eq-flatsize}.  Cluster energies are disjoint and enlarged
physical time bins have bounded overlap, so this is a global sum.
\end{proof}

\subsection{Density bands and the frame estimate}

The two budgets will serve different parts of the summation.  Very
small densities will pay the polynomial number of clusters and time
bins, so their energies can be bounded one time bin at a time.  For
the remaining densities we will choose transverse coarse bins from
one common density band.  The corresponding density weight can then
remain attached to each packet energy and be summed with the global
weighted budget.  We now construct these two kinds of tuples and verify the
hypotheses of the frame estimate.

Write the densities in dyadic bands $d_\ell=C_d2^{-\ell}$,
$\ell\geq0$, with fixed upper constant $C_d$.  Only $O(\log h)$
indices occur on the needed cells.  Indeed their distances from the
measure support are bounded by a fixed power of $h$, while
\eqref{w:eq-incidence-mass} gives a polynomial lower bound on total
mass; the strictly positive kernel in \eqref{w:eq-flatdensity} then
gives $d_{\theta,T}\geq h^{-C}$.

At a witness for each significant coarse bin, at least one band field
has modulus at least
\begin{equation}\label{w:eq-bandlevel}
 A_\ell=c'R^{-n}\lambda(1+\ell)^{-2}.
\end{equation}
This follows by threshold splitting, since
$\sum_{\ell\geq0}(1+\ell)^{-2}<\infty$.  Choose one such band for
each significant bin.  There are two possibilities:
\begin{enumerate}[label=\textup{(\roman*)}]
\item A chosen band is tiny, $d_\ell<X^{-C_{11}}$.  Starting from its
center, complete it to a robust $n=k+1$ tuple of significant coarse
bins in tangential frequency space.
\item All chosen bands are non-tiny.  There are only $O(C_{11}\log X)$
colors.  Some single color supplies a robust $n$-tuple.
\end{enumerate}
Both statements follow from the strip-union exclusion: successive
maximal affine heights can all be chosen at least $C_*/R$.  If a color
failed to supply such a simplex, it would lie in one strip.  The
determinant lower bound for a successful tuple is $cR^{-k}$ throughout
the bins when $C_*$ is large enough.

Fix a cluster, a time bin of length $h^2$, an ordered coarse tuple and
its band indices $\ell_i$, $1\leq i\leq n$.  Let $M_0$ be the mass
of tested points assigned to this tuple, and let
$E_i=\sum s_{\theta,T}^2$ over its $i$th part, with the fixed time
enlargement.  A tuple of mass at most $mh^{2n-c_4}$ can be discarded,
where $c_4>0$ is a small fixed dimensional constant.  The total number
of tuples is $X^C(\log h)^C$, so their relative discarded mass, by
\eqref{w:eq-incidence-mass}, is at most
\[
 X^C(\log h)^C(h/L)^{2n}h^{-c_4}=o(1).
\]
This is one of the requirements imposed on the later $\kappa$.

For every remaining tuple, Markov and \eqref{w:eq-flatdensity} retain
at least half its mass where all full square envelopes are bounded by
\begin{equation}\label{w:eq-Bi}
 B_i=C\left(\frac{mh^{n+1}d_{\ell_i}E_i}{M_0}\right)^{1/2}.
\end{equation}
They remain comparable at every point of the macro containing the
tested point, since $R^4\ll h$.  Thus they also bound the envelopes
at the separate coarse-bin witnesses.  Zero-energy cases need no
consideration.  By \eqref{w:eq-flat-unweighted} and the mass cutoff,
\begin{equation}\label{w:eq-framecutoff}
 B_i\lesssim h^{(1-2n+c_4)/2},\qquad A_{\ell_i}/B_i\geq h^{c_3}
\end{equation}
for a fixed $c_3>0$.  For the latter assertion, use
$g\geq D^{-1}$, $\chi>L^{-v_0}$ and the margin
\eqref{w:eq-threshold-margin}.  Choose $c_4,v_0,c_3$ as sufficiently
small dimensional fractions of that margin; the later choice of
$\kappa$ pays every fixed $X$ power in $L/h$, $R$ and the threshold,
and logarithmic factors are smaller than any fixed positive power of
$h$.  This verifies the frame cutoff without allowing a geometric
power to determine $u$ retroactively.

\begin{lemma}[A tuple bound]\label{w:tuple}
For each retained tuple,
\begin{equation}\label{w:eq-tuple}
 \lambda^2 M_0\leq CR^C(1+\textstyle\sum_i\ell_i)^C
 m\chi^{1-\beta}h^{n+2n/(n+1)}
 \left(\prod_{i=1}^n d_{\ell_i}E_i\right)^{\beta/n}
 \left(\sum_{i=1}^nE_i\right)^{1-\beta}.
\end{equation}
The $R$ exponent is independent of the number of plate clusters and
time bins and of $C_{10},C_{11}$.
\end{lemma}

\begin{proof}
Partition the sheared normal spatial coordinate into slabs of width
$h$, with bounded enlargements to accommodate witnesses.  The projected
measure in tangential spacetime has unit-cube mass at most $Cm\chi h$
on each slab, by unit covering in the normal coordinate.

Only nearby normal and time labels of compact packets contribute in
a fixed slab and time bin.  Recenter at the time-bin origin.  The
residual normal frequency is $O(1/h)$, so its drift during $h^2$ is
$O(h)$.  Projected tangential packets consequently have bounded
phase-space multiplicity.  Rotating the lattice changes only bounded
multiplicities and comparable polynomial weights.  The omitted normal
and time indices have bounded normalized displacements, so the projected
envelopes are bounded by a constant times the full envelopes.

In the normal coordinate divided by $h$, expand each cutoff profile
in a Fourier series.  The residual normal carrier and time phase have
bounded normalized derivatives and can be included in that profile.
The exterior Fourier phase has modulus one.  For mode $m\in\Z$,
the tangential packet amplitude and envelope gain
$\rho_m\lesssim(1+|m|)^{-A_*}$, with any required fixed $A_*$ supplied
by the original derivative order.  Split thresholds by
$\gamma_m=c(1+|m|)^{-2}$.  The frame cutoff improves rather than
deteriorates under the factor $\gamma_m/\rho_m$.

Witnesses need not be on the tested point's projected unit cube.
Translate the tangential spacetime packets by an integer offset,
separately for each family, so that its witness lies in that cube.
Each offset is $O(R^4)$, so summing all choices costs $R^C$.
Counting, transversality and envelope comparability are preserved.

Apply \Cref{f:frame}, normalizing family $i$ by $B_i$, and multiply
the resulting cube count by $Cm\chi h$.  Before summing modes, the
right side contains the decay factors
\[
 \prod_{i=1}^n(\rho_{m_i}/\gamma_{m_i})^{4/(kn)}
 \quad\hbox{and, in its $j$th summand,}
 \quad(\rho_{m_j}/\gamma_{m_j})^2.
\]
Every family has a positive decay exponent, even when it does not
index the energy summand.  All independent mode sums therefore
converge for fixed large $A_*$.  Sum slabs at this stage, keeping the
same global $B_i$ from \eqref{w:eq-Bi} in every slab and putting only
its nearby packet energies in the linear sum.  Each packet occurs in
boundedly many enlarged slabs.  We obtain
\begin{equation}\label{w:eq-frame-before-solve}
 M_0\leq CR^C m\chi h
 \prod_{i=1}^n(B_i/A_{\ell_i})^{4/(kn)}
 \sum_{i=1}^n\frac{h^{k+2}E_i}{A_{\ell_i}^2}.
\end{equation}
Summing slabs after solving a nonlinear mass inequality instead would
lose a multiplicity factor; \eqref{w:eq-frame-before-solve} avoids it.

Substitute \eqref{w:eq-Bi} and \eqref{w:eq-bandlevel}.  Up to the
displayed $R$ and band-index factors, the result is
\[
 M_0^{1+2/k}\leq C R^C(1+\textstyle\sum_i\ell_i)^C
 m^{1+2/k}\chi
 h^{n+2+2(n+1)/k}\lambda^{-2-4/k}
 \left(\prod_i d_{\ell_i}E_i\right)^{2/(kn)}\sum_iE_i.
\]
Raise both sides to $k/(k+2)=1-\beta$.  The exponent of $h$ becomes
$n+2n/(n+1)$ and the product exponent becomes $\beta/n$; the
$\lambda$ exponent is exactly $-2$.  This is
\eqref{w:eq-tuple}.
\end{proof}

\subsection{Summation and completion of the induction}

\begin{proof}[Completion of the proof of \Cref{w:weak}]
First sum the tuples with a tiny band.  The per-bin bound
$E_i\lesssim h^{-n}$ gives
\[
 \left(\prod_i E_i\right)^{\beta/n}
 \left(\sum_iE_i\right)^{1-\beta}\lesssim h^{-n}.
\]
Here the polynomial number of clusters and the
$O((L/h)^2)=X^{2C_{10}}$ time bins may be counted explicitly.  Their
total power of $X$ is fixed before $C_{11}$.  To sum the density bands,
extend the range to all nonnegative indices and use
\begin{equation}\label{w:eq-tinysum}
 \sum_{\min_i d_{\ell_i}<X^{-C_{11}}}
 (1+\textstyle\sum_i\ell_i)^C\prod_i d_{\ell_i}^{\beta/n}
 \leq C X^{-\beta C_{11}/(2n)}.
\end{equation}
Indeed assign a tiny index, bound the polynomial by a product of
polynomials, and absorb each into half its exponential decay.  The
remaining series are geometric.  There is no factor $\log h$ in
\eqref{w:eq-tinysum}.  Choose $C_{11}$ large enough to pay all the
previous cluster, time-bin and tuple powers and the desired negative
power of $D$.

For a common non-tiny band, the energy expression in
\eqref{w:eq-tuple} is bounded by
\[
 C d_\ell^\beta\sum_i E_i.
\]
Each packet has one density index, and can occur in only $R^C$ coarse
tuple and witness-offset choices.  Normal slabs and enlarged time
bins have bounded overlap.  Cluster labels are disjoint coefficient
subarrays.  Hence this sum is charged directly to the global weighted
bound \eqref{w:eq-flat-weighted}, with no factor counting clusters or
time bins.  Also
$\ell\leq C+C_{11}\log_2X$, so its band factor is only
$C(C_{11})(1+\log X)^C$.  The scale factors cancel as
\[
 h^{n+2n/(n+1)}h^{-n}b_1^{2n/(n+1)}
 =L^{2n/(n+1)}.
\]
The choices of $u$ and then $B$ absorb $R^{C_f}$ and the fixed
logarithmic powers into part of the sparse gain.  Combining both
types of tuples, and recalling that a fixed fraction of $M_Y$ remained,
gives
\begin{equation}\label{w:eq-finalL2}
 \lambda^2M_Y\leq C_Bm\chi^{1-\beta}
                   L^{2n/(n+1)}D^{-c_s/2}.
\end{equation}

We now verify the promised polynomial level bound, before choosing
$p-2$.  From \eqref{w:eq-incidence-mass},
$M_Y\geq cmL^{2n}X^{-C_{\mathrm{geo}}}$.  Substitute this and
$\lambda=\chi^agL^{-e}$ in \eqref{w:eq-finalL2}.  Since
$2n-2e=2n/(n+1)$, we obtain
\[
 g^2\leq C_B\chi^{\zeta}X^{C_{\mathrm{geo}}}D^{-c_s/2}
 \leq C_BX^{C_{\mathrm{geo}}}.
\]
Thus $g\leq C_BX^{C_g}$ with $C_g$ fixed by geometry, independently
of $\kappa,p-2,C_{\mathrm{ind}}$.  This confirms the noncircular order
of choices in \Cref{w:choices}.

Finally \eqref{w:eq-finalL2} and $\chi^{\zeta}\leq1$ imply
\[
 M_Y\leq C_BmL^{2n}g^{-2}D^{-c_s/2}
 \leq C'_BmL^{2n}g^{-p}D^{-\sigma},
\]
provided
\[
 C_g(p-2)+\sigma<c_s/2.
\]
This requirement is compatible with the easy-range inequalities,
since $p-2$ is chosen after all geometric exponents and after $\kappa$.
All uses of induction removed small fractions of the hypothetical
mass: their ratios were negative powers of $X$, or
$R^{-c_2}D^\sigma$, fixed before choosing $C_{\mathrm{ind}}$.
Every remaining constant $C'_B$ is independent of that induction
constant.  Taking $C_{\mathrm{ind}}>C'_B$ contradicts
\eqref{w:eq-counterlevel}.  Both inductions close, proving
\Cref{w:weak}.
\end{proof}
